\chapter{Objectius}
Quan vaig començar a plantejar aquest treball de recerca, tenia molt clar que volia que tingués una part pràctica relacionada amb la programació i les dades reals. A partir d’aquí vaig definir una sèrie d’objectius que m’han servit per orientar tot el procés i que explico a continuació.

\section{Objectiu general}
L’objectiu general del meu treball és \textbf{crear un programa en Python que utilitzi dades reals del servei Bicing de Barcelona per calcular els canvis ràpids possibles a cada estació}. Amb això no només vull aprendre a programar amb informació real sinó també entendre millor com funciona un sistema de mobilitat compartida i quines dificultats pot tenir en el dia a dia.

\section{Objectius específics}
Per arribar a l’objectiu principal, he definit diversos objectius més concrets:
\begin{itemize}
    \item Aprendre a utilitzar una API: entendre què és, com funciona i com es poden obtenir dades en temps real d’un servei públic com el Bicing.
    \item Treballar amb el format JSON: aprendre a llegir i processar dades en aquest format, que és el més habitual en aplicacions modernes.
    \item Programar amb Python: posar en pràctica coneixements de Python per llegir dades, recórrer-les i aplicar càlculs senzills.
    \item Definir i calcular el concepte de ``canvi ràpid'': establir una mesura senzilla però útil que mostri en quines estacions hi ha tant bicicletes com espais disponibles.
    \item Mostrar els resultats de manera clara: dissenyar una sortida de dades entenedora, amb informació de cada estació i els canvis ràpids possibles.
    \item Analitzar els resultats obtinguts: reflexionar sobre el que ens diuen les dades, veure quines estacions estan equilibrades i quines no, i pensar com això afecta els usuaris.
    \item Fer una memòria escrita completa: redactar tot el procés amb introducció, marc teòric, metodologia, resultats i conclusions, seguint el format d’un treball de recerca.
\end{itemize}

\section{Objectius personals}
A més dels objectius estrictament tècnics, també m’he proposat objectius personals:
\begin{itemize}
    \item Millorar la meva organització i capacitat de treball a llarg termini.
    \item Aprendre a resoldre problemes reals (errors de codi, connexió a la xarxa, interpretació de dades).
    \item Apropar-me al món de la programació i de les dades obertes, que pot ser útil en estudis futurs.
\end{itemize}

\section{Preguntes de recerca}
Per guiar encara més el treball, també m’he plantejat algunes preguntes concretes que intentaré respondre amb el programa:
\begin{itemize}
    \item Quines estacions de Bicing permeten més canvis ràpids en un moment determinat?
    \item Hi ha diferències entre estacions cèntriques i perifèriques?
    \item Què passa quan una estació té moltes bicis però pocs espais, o a l’inrevés?
    \item Podria aquest mètode ajudar a gestionar millor el servei?
\end{itemize}
